\section{Parse Tree dan Hubungannya dengan AST}
\label{sec:ch08-ast}

\subsection{Parse Tree}

Pohon urai (\eng{parse tree}) merepresentasikan derivasi secara
lengkap. Guntoro menyebutnya pohon sintaks atau pohon penurunan
(\eng{derivation tree}). Setiap simpul dalam adalah nonterminal,
setiap daun adalah terminal atau \(\varepsilon\), dan anak-anak
simpul sesuai dengan ruas kanan aturan yang dipakai.

Pohon urai memuat semua detail sintaksis, termasuk tanda kurung,
titik koma, dan kata kunci. Detail ini penting untuk validasi, tetapi
membuat pohon besar dan sulit diproses pada fase berikutnya.

\subsection{Abstract Syntax Tree}

Abstract Syntax Tree (AST) adalah penyederhanaan pohon urai. Juarna
menyebutnya \emph{syntax tree}. AST menghilangkan detail sintaksis
yang tidak diperlukan untuk analisis semantik, seperti tanda kurung
dan titik koma. Setiap simpul AST merepresentasikan konstruksi bahasa:
ekspresi, pernyataan, deklarasi, atau fungsi.

\begin{figure}[H]
\centering
\begin{minipage}{0.46\textwidth}
\centering
\begin{forest}
parsetree
[E
  [E [T [F [\code{id}]]]]
  [+]
  [T
    [T [F [\code{id}]]]
    [*]
    [F [\code{id}]]
  ]
]
\end{forest}
\end{minipage}\hfill
\begin{minipage}{0.46\textwidth}
\centering
\begin{forest}
asttree
[+
  [\code{id}]
  [*
    [\code{id}]
    [\code{id}]
  ]
]
\end{forest}
\end{minipage}
\caption{Pohon urai (kiri) dan AST (kanan) untuk ekspresi yang sama}
\label{fig:ch08-parse-vs-ast}
\end{figure}

\subsection{Perbedaan Penting}

\begin{table}[H]
\centering
\caption{Perbedaan pohon urai dan AST}
\label{tab:ch08-parse-ast}
\begin{tabularx}{\textwidth}{@{}lXX@{}}
\toprule
\textbf{Aspek} & \textbf{Pohon urai} & \textbf{AST} \\
\midrule
Isi & Semua simbol tata bahasa & Konstruksi bermakna \\
Ukuran & Lebih besar & Lebih ringkas \\
Detail & Kurung, titik koma, kata kunci & Hanya operand dan operator \\
Penggunaan & Analisis sintaksis & Analisis semantik dan kode \\
\bottomrule
\end{tabularx}
\end{table}

\subsection{Contoh AST}

Untuk program \file{arith.sc}, AST tingkat atas memuat deklarasi
fungsi \code{main}, deklarasi variabel \code{a} dan \code{b}, serta
pernyataan \code{print}. Ekspresi \code{2 + 3 * 4} menjadi AST
dengan akar \code{+}, anak kiri \code{2}, dan anak kanan \code{*}
dengan anak \code{3} dan \code{4}. Struktur ini langsung
mencerminkan prioritas operator.

\subsection{Dari Pohon Urai ke AST}

Konversi dapat dilakukan saat parsing. Ketika parser mengurangi
(\eng{reduce}) aturan produksi, ia membangun simpul AST yang sesuai.
Aturan yang hanya mengatur sintaksis, seperti
\(F \to (E)\), tidak menghasilkan simpul baru; ia hanya meneruskan
simpul \(E\). Aturan operator menghasilkan simpul operator dengan
anak operand.

Pendekatan ini menghindari pembentukan pohon urai lengkap, sehingga
lebih efisien. Namun, untuk keperluan pengajaran, pohon urai tetap
digambar agar mahasiswa memahami hubungan antara derivasi dan AST.
